\section{Fixed-parameter decision and radius}\label{sec:fpt}

\subsection{Bounded dynamic programs}

Fix a pose \(A\) of the canonical largest piece \(L\), and retain the
notation \(b=X(T)-X(A)\).  For margin feasibility, a partial-placement
dynamic program needs only residual margin states bounded coordinatewise by
\(b\).  Its exact state box has size
\[
  B(b)=\prod_j(b_j+1).
\]
For nonnegative integers \(n\), \(n+1\leq2^n\).  The row coordinates of
\(b\) sum to \(q\), as do the column coordinates, so
\[
  B(b)\leq 2^q2^q=4^q.
\]
For exact coverage, the residual target \(T\setminus A\) has \(q\) cells.
A subset-mask dynamic program therefore has at most \(2^q\) states.
Disjoining over the poses of \(L\) decides whether the margin fiber and the
exact-cover sublocus are empty.

To decide purity and repair, we reconstruct all accepting margin paths.  This
is safe because the entire fiber is itself parameter-bounded.

\begin{theorem}[Largest-piece residual-area FPT theorem]\label{thm:fpt}
For an area-compatible fixed selected set of named finite grid pieces and a
finite orientation group of size at most \(h\), support-exactness, fiber
purity, and repair-exactness in the exact-two graph are fixed-parameter
tractable in \((q,h)\).  The full margin fiber satisfies
\[
 |\FFib_E|\leq h(q+1)(h q^2)^q
              =2^{O(q\log(qh))}.
\]
For fixed \(h\), this simplifies to \(2^{O(q\log q)}\).  The replacement
graph can be constructed in
\[
 2^{O(q\log(qh))}\operatorname{poly}(|I|)
\]
time, and hence in \(2^{O(q\log q)}\operatorname{poly}(|I|)\) time when
\(h\) is fixed.
\end{theorem}

\begin{proof}
By \cref{cor:pose-bounds}, \(L\) has at most \(h(q+1)\) legal poses.  For a
fixed largest pose, there are at most \(q\) residual pieces and each has at
most \(hq^2\) margin-compatible candidate poses.  Hence
\[
 |\FFib_E|
 \leq h(q+1)(hq^2)^q.
\]
This includes nonaccepting combinations in its upper estimate and is
therefore valid for the fiber.  The case \(q=0\) is handled directly.

Enumerating the parameter-bounded fiber and testing the cell coverage of
each tuple decides purity.  Support is decided by comparing nonemptiness of
\(\FFib_E\) and \(\TT_E\).

Hash all fiber tuples.  For each tuple, choose two named coordinates,
enumerate alternatives from their candidate-pose catalogues, and retain
the resulting tuple precisely when it occurs in the hash table.  There are
at most polynomially many choices outside the parameter-bounded pose
factors, so this constructs \(\GG_E^{(2)}\) in
\(2^{O(q\log(qh))}\operatorname{poly}(|I|)\) time.  Mark its connected
components and test whether each intersects \(\TT_E\).  This decides
repair-exactness.
\end{proof}

\begin{remark}
Adding the one-name cases constructs \(\GG_E^{(\leq2)}\) in the same FPT
class.  This does not identify the two graphs or transfer finite component
counts between them.
\end{remark}

\begin{corollary}[Existence bound for repair radius]\label{cor:radius}
Every repair-exact fiber satisfies
\[
 \rho(E)\leq|\FFib_E|-1
 \leq h(q+1)(h q^2)^q-1.
\]
\end{corollary}

\begin{proof}
In a component meeting \(\TT_E\), choose a shortest path from a vertex to a
tiling.  A shortest path is simple and hence has at most the component size
minus one edges.
\end{proof}

The bound in \cref{cor:radius} is an existence bound.  It is not asserted to
be polynomial, linear, sharp, or representative of observed runtime.
